\documentclass[11pt]{article}
\usepackage[T1]{fontenc}
\usepackage{amsmath,amssymb,amsthm}
\usepackage[margin=1in]{geometry}

\newtheorem{theorem}{Theorem}
\newtheorem{lemma}{Lemma}
\newtheorem{corollary}{Corollary}
\newcommand{\NBC}{\operatorname{NBC}}
\newcommand{\Vol}{\operatorname{Vol}}
\newcommand{\calT}{\mathcal{T}}
\newcommand{\calO}{\mathcal{O}}
\newcommand{\1}{\mathbf{1}}

\title{A Human-Readable Account of the NBC Volume Identity for Hard-Sphere Mayer Sums}
\author{ukiyois (research direction and maintainer)\\
AI-assisted Lean formalization and proof extraction}
\date{2026-08-27}

\begin{document}
\maketitle

\begin{abstract}
For a hard-sphere configuration, the connected Mayer graph sum is a signed
connected-subgraph count of its overlap graph.  A finite broken-circuit
cancellation turns that count into a fixed-sign number of no-broken-circuit
spanning trees.  Regrouping over the corresponding tree-owned regions gives
the exact positive volume identity.  The final statement is formulated on the
standard anchored Euclidean configuration space with product Lebesgue measure.
An explicit measure-preserving coordinate map to the flat space
$\mathbb{R}^{d(k-1)}$ is also recorded, so the two common presentations of the
configuration measure are not silently identified.  The formal companion
checks both the finite and analytic steps.
\end{abstract}

This document is an explanatory mathematical translation of the formalization:
it records the proof architecture and the key intermediate identities in a
reader-friendly form.  The kernel-checked proof is the Lean development in
this repository; this TeX file is not itself a formal proof artifact.

\section{Definitions}

Fix $k\ge2$, dimension $d\ge1$, diameter $1$, and $r_1=0$.  For a
configuration $x=(r_2,\ldots,r_k)\in\mathbb{R}^{d(k-1)}$, define its overlap
graph
\[
 E(\calO(x))=\{\{i,j\}:|r_i-r_j|<1\}.
\]
For a graph $\Gamma$ on $[k]$, let $c_\Gamma(A)$ be the number of connected
components of $([k],A)$, and define
\[
 m(\Gamma)=\sum_{\substack{A\subseteq E(\Gamma)\\c_\Gamma(A)=1}}
              (-1)^{|A|}.
\]
Fix a strict edge order.  A broken circuit is a circuit with its largest edge
removed.  Write $\NBC(\Gamma)$ for the number of spanning trees containing no
broken circuit.

Let $\calT_k$ denote the finite set of spanning trees of the complete graph
$K_k$.  Equip the anchored configuration space with its product Lebesgue
measure; the notation $dx$ below refers to this canonical measure.

For completeness, there is an explicit bridge to the flat-coordinate model
used in much of the virial literature.  Write

\[
 X=\operatorname{Fin}(k-1)\to\operatorname{EuclideanSpace}(\mathbb{R},
     \operatorname{Fin}(d)),\qquad
 Y=\operatorname{EuclideanSpace}(\mathbb{R},\operatorname{Fin}((k-1)d)).
\]

The map $F:X\to Y$ first writes each Euclidean block in its real
coordinates, then sends $(i,j)$ to the corresponding coordinate of
\(\operatorname{Fin}((k-1)d)\) using the standard finite-index equivalence.
Thus $F$ is a bijective coordinate reordering.  The measure proof factors
into four elementary steps: the Euclidean coordinate representation in each
block, finite-product currying, reindexing by the finite-index equivalence,
and the inverse Euclidean coordinate representation.  Each step preserves
the corresponding product volume, so

\[
 F_*\,dx=dy.
\tag{0}
\]

In particular, for every measurable $A\subseteq X$,

\[
 \Vol_Y(F(A))=\Vol_X(A).
\tag{0a}
\]

The formal development defines the flat NBC region as $F(\NBC_T)$ and proves
(0a) for it.  This makes the usual $\mathbb{R}^{d(k-1)}$ presentation an
explicit transported version of the product-space statement, without
introducing a second unproved measure convention.

For a spanning tree $T$ of $K_k$, put
\[
 X_T=\{x:|r_i-r_j|<1\text{ for every }ij\in E(T)\},
\]
and
\[
 \NBC_T=\{x\in X_T:T\text{ contains no broken circuit of }\calO(x)\}.
\]
The definition of $X_T$ constrains tree edges only; other pairs may also be
active.

The connected Mayer weight and anchored cluster integral are
\[
 \omega_k(x)=\sum_{G\text{ connected on }[k]}
       \prod_{e\in E(G)}f_e(x),
 \qquad f_e(x)=-\1_{e\in E(\calO(x))},
\]
\[
 b_k=\frac{1}{k!}\int\omega_k(x)\,dx.
\]

\section{Restriction to the active graph}

\begin{lemma}
For every configuration $x$,
\[
 \omega_k(x)=m(\calO(x)).
\]
\end{lemma}

\begin{proof}
For a connected graph $G$, its product is zero unless
$E(G)\subseteq E(\calO(x))$.  In the latter case the product is
$(-1)^{|E(G)|}$.  Replacing $E(G)$ by $A$ gives exactly the sum defining
$m(\calO(x))$.
\end{proof}

\section{The finite broken-circuit cancellation}

\begin{lemma}[NBC cancellation]
For a connected graph $\Gamma$ on $[k]$,
\[
 m(\Gamma)=(-1)^{k-1}\NBC(\Gamma).
\tag{1}
\]
\end{lemma}

\begin{proof}
Call an edge $e$ a candidate for an edge set $A$ if there is a circuit $C$
whose largest edge is $e$ and such that $C\setminus\{e\}\subseteq A$.  Thus
$A$ contains a broken circuit exactly when it has a candidate.  For a bad
edge set $A$, choose the least candidate $e(A)$ and define
\[
 \tau(A)=A\mathbin{\triangle}\{e(A)\}.
\]
Let $C$ witness that $e(A)$ is a candidate.  The path
$C\setminus\{e(A)\}$ joins the endpoints of $e(A)$, so adding or removing
$e(A)$ does not change the connected components of $([k],A)$.  The same broken
circuit remains contained after the toggle, so $e(A)$ remains a candidate.

If $f<e(A)$, then $e(A)$ cannot belong to a broken circuit whose largest edge
is $f$, since every edge in that broken circuit is smaller than $f$.  Hence
the candidate status of every smaller edge is unchanged by the toggle.  It
follows that $e(\tau(A))=e(A)$ and $\tau(\tau(A))=A$.  The map $\tau$ is
therefore a fixed-point-free involution on the bad connected edge sets, and it
changes their cardinality by one.  Their signs cancel in pairs.

It remains to identify the fixed terms.  A connected edge set containing a
cycle contains the broken circuit obtained by deleting the largest edge of
that cycle.  Thus a connected edge set with no broken circuit is acyclic and
connected, hence is a spanning tree.  Conversely every NBC spanning tree is
not bad.  The only surviving terms in $m(\Gamma)$ are consequently the
$\NBC(\Gamma)$ trees, each with $k-1$ edges, proving (1).
\end{proof}

\begin{lemma}[Tutte specialization]
For connected $\Gamma$,
\[
 T_\Gamma(1,0)=\NBC(\Gamma),\qquad
 m(\Gamma)=(-1)^{k-1}T_\Gamma(1,0),
\]
where
\[
 T_\Gamma(u,v)=\sum_{A\subseteq E(\Gamma)}
 (u-1)^{c_\Gamma(A)-1}(v-1)^{|A|-k+c_\Gamma(A)}.
\]
\end{lemma}

\begin{proof}
At $u=1$, only $c_\Gamma(A)=1$ survives.  At $v=0$ the remaining term is
$(-1)^{|A|-k+1}$, which is $(-1)^{k-1}$ times the sign in $m(\Gamma)$.
Equation (1) then gives $T_\Gamma(1,0)=\NBC(\Gamma)$.
\end{proof}

\section{Pointwise NBC identity}

From the two preceding lemmas, for every configuration $x$,
\[
 \omega_k(x)=
 \begin{cases}
 (-1)^{k-1}\NBC(\calO(x)),&\calO(x)\text{ connected},\\
 0,&\calO(x)\text{ disconnected}.
 \end{cases}
\tag{2}
\]
Let
\[
 n(x)=\#\{T\in\calT_k:x\in\NBC_T\}.
\]
If $\calO(x)$ is connected, $x\in X_T$ is equivalent to
$E(T)\subseteq E(\calO(x))$, and the second condition in $\NBC_T$ is exactly
the NBC condition in $\calO(x)$.  If $\calO(x)$ is disconnected, neither side
counts a spanning tree.  Thus
\[
 |\omega_k(x)|=n(x)=\sum_{T\in\calT_k}\1_{\NBC_T}(x).
\tag{3}
\]

\section{NBC volume identity}

The analytic details needed to use (3) are elementary but must be included.
Every edge activity is the indicator of an open norm inequality, so every
$X_T$ and $\NBC_T$ is Borel: there are only finitely many possible active
edge patterns, and the NBC condition is a finite Boolean formula in those
activities.  If $\calO(x)$ is connected, a path from $1$ to any vertex has at
most $k-1$ edges, each of length less than $1$.  Hence every coordinate of
$x$ has norm less than $k-1$.  If $\calO(x)$ is disconnected, $\omega_k(x)=0$.
Thus the integrand is bounded and supported in a finite-radius closed ball;
it is integrable, and every $\NBC_T$ has finite volume.

\begin{theorem}
For every $k\ge2$,
\[
 k!\,|b_k|=\int |\omega_k(x)|\,dx
 =\sum_{T\in\calT_k}\Vol(\NBC_T).
\tag{4}
\]
\end{theorem}

\begin{proof}
Equation (2) gives a constant-sign integrand, so
$|\int\omega_k|=\int|\omega_k|$.  Substitute (3).  The tree set
$\calT_k$ is finite, and all indicators are measurable and nonnegative, so
finite additivity of the integral gives
\[
 \int\sum_{T\in\calT_k}\1_{\NBC_T}(x)\,dx
 =\sum_{T\in\calT_k}\int\1_{\NBC_T}(x)\,dx
 =\sum_{T\in\calT_k}\Vol(\NBC_T).
\]
Finally use $b_k=k!^{-1}\int\omega_k$.
\end{proof}

\section{Provenance and scope}

The research direction supplied by \texttt{ukiyois} was the key graph-theoretic
and NBC approach to the signed volume problem, together with the question
whether a positive decomposition might exist; \texttt{ukiyois} did not state
the exact identity proved here.  The upstream verified fact records are
identified by their internal worker authors as follows: fact
\texttt{2041d4f28d05949e} by \texttt{xhigh3}, fact
\texttt{4afee99f8b8e6010} by \texttt{xhigh5}, fact
\texttt{6d35ccd3e0232041} by \texttt{high3}, and fact
\texttt{7dcc52cfec5875b2} by \texttt{xhigh5}.  They are provenance records
only and are not imported as proof terms.  The Lean definitions and proofs
were implemented by AI agent sessions, principally DeepSeek-V4-Flash and
GLM-5.3-Flash through OpenCode; some runs exposed only the anonymous provider
alias \texttt{ox-alpha}.  The coordinate measure bridge and this final
standalone extraction were added in an OpenCode session using
\texttt{openai/gpt-5.6-luna}.  No claim of strict positivity of $b_k$ or of a
separate complete sign-alternation theorem is made here.

\section{Formal verification boundary}

The Lean companion does not import fact statements as proof terms and adds no
project-specific axiom.  The files
\texttt{NBCMatroid.lean}, \texttt{NBCGraph.lean}, and
\texttt{GraphicMatroid.lean} check the finite cancellation and graph
characterizations, while \texttt{MayerNBC.lean} checks the abstract Mayer
restriction, the pointwise NBC decomposition, and the finite-measure region
sum.  The dedicated file \texttt{HardSphereNBC.lean} specializes this chain
to
\[
  \texttt{HardSphereConfiguration k d}
    := \texttt{Fin (k - 1) -> HSPosition d},
\]
with particle $1$ at the origin and the product of Euclidean volume measures.
It proves the Borel measurability, bounded-support integrability, real-valued
absolute-value identity, and theorem (4) as
\texttt{hardSphere\_nbc\_volume\_identity}.  Fact IDs are provenance only;
they are not imported as proof terms.  The companion
\texttt{HardSphereMeasure.lean} proves (0), (0a), and the flat-coordinate form
\texttt{hardSphere\_nbc\_volume\_identity\_flat}.

\end{document}
